\paragraph{Reading.} LUDVIG's curves coincide with ours at every floor, as expected from identical tables. Occam's
curves cross ours. Its peaks are slightly higher, so at a given floor it keeps a few more candidates. When the methods
are compared at an equal number of absent objects firing, which is the fair comparison:
\begin{center}\small
\begin{tabular}{lcccc}
\toprule
Absent objects with a candidate (of 6) & 1 & 2 & 3 & 4\\
\midrule
Ours / LUDVIG, Old & 19 & \textbf{23} & 23 & 23\\
Ours / LUDVIG, A   & 19 & 22 & 23 & 23\\
Occam's, Old       & 19 & 21 & \textbf{24} & \textbf{24}\\
Occam's, A         & 19 & 21 & \textbf{24} & \textbf{24}\\
\bottomrule
\end{tabular}
\end{center}
\noindent (best hits over the floors 0.46--0.60 with at most that many absent objects firing). Occam's is ahead by one query
at 3--4 false positives and behind by one or two at 2; at 1 false positive every method finds 19. The ceiling over
all floors is 24/29 for every lift and teacher set, the same ceiling the 8 Oct sweep found for B2. As on 8 Oct, the
same two absent objects keep firing whatever the lift: \emph{ceiling fan} (backroom) at every floor for every table,
and \emph{recycling bin} (flightroom) up to floor 0.55 with the Old teachers and 0.575 with A.
These floors are read off the evaluation set and are optimistic.
